\begin{figure}[H]
\centering
\makebox[\textwidth][c]{%
\begin{tikzpicture}[
    scale=0.9,
    transform shape,
    actor/.style={circle, draw=black, fill=white, line width=0.8pt, minimum size=0.7cm},
    usecase/.style={ellipse, draw=black, fill=white, line width=0.55pt, minimum width=3.4cm, minimum height=0.7cm, font=\scriptsize, align=center},
    system/.style={rectangle, draw=black!50, dashed, fill=white, line width=0.4pt, inner sep=8pt, rounded corners=3pt}
]

\node[system, minimum width=12cm, minimum height=7cm] (system) at (0,0) {};
\node[font=\bfseries\small] at (0,3.1) {Phân rã: Nhập kho \& Xuất kho};

\node[usecase] (uc22) at (-3.0,1.3) {UC22: Lập phiếu nhập};
\node[usecase] (uc23) at (3.0,1.3) {UC23: Duyệt phiếu nhập};
\node[usecase] (uc24) at (-3.0,-0.5) {UC24: Lập phiếu xuất};
\node[usecase] (uc25) at (3.0,-0.5) {UC25: Duyệt phiếu xuất};

\node[actor, label=below:{\scriptsize\bfseries Nhân viên kho}] (nvk) at (-6.5,0.4) {};
\node[actor, label=below:{\scriptsize\bfseries Quản lý kho}] (qlk) at (6.5,0.4) {};

\draw (nvk) -- (uc22);
\draw (nvk) -- (uc24);
\draw (qlk) -- (uc23);
\draw (qlk) -- (uc25);

\end{tikzpicture}
}
\caption{Sơ đồ use case phân rã -- Nhập/xuất (lập và duyệt là 2 UC khác nhau)}
\label{fig:usecase-inventory-management}
\end{figure}
